\section{Convolutional Neural Networks (6 Points)}

\begin{enumerate}

\item \textbf{[2pts]} \textbf{Numerical answer:} Suppose the first layer of a CNN uses a $5 \times 5$ convolutional filter and outputs 3 channels. It takes as input images of size $100 \times 100$. It also uses a padding of $3$ and a stride $2$ in both input dimensions. What are the dimensions of the output of this layer?

\begin{tcolorbox}[fit,height=1cm, width=8cm, blank, borderline={1pt}{-2pt},nobeforeafter]
%solution
\end{tcolorbox}

\item \textbf{[2pts]} \textbf{Numerical answer:} Suppose you have a CNN that takes inputs of shape $3 \times 100 \times 100$. Suppose the first convolutional layer uses $3\times 5 \times 5$ filters and outputs $8$ channels. How many parameters does this layer have (including a scalar bias for each output channel)?

\begin{tcolorbox}[fit,height=1cm, width=8cm, blank, borderline={1pt}{-2pt},nobeforeafter]
%solution
\end{tcolorbox}


\item \textbf{[2 pts]} \textbf{Select all that Apply:} Which of the following statements about pointwise (sometimes called $1\times 1$) convolutions is/are true?
\begin{quote}
    $\Box$ Pointwise convolutions introduce non-linearity between two CNN layers \newline
    $\Box$ A pointwise convolution can be formulated as a fully connected layer that maps each stack of pixels (across all channels) in the input to a pixel in the output
    \newline
    $\Box$ Pointwise convolutions can reduce the overall number of parameters in a model by reducing the total size of the following layer 
    \newline
    $\Box$ Pointwise convolutions preserve the number of channels but may impact other dimensions \newline
    $\Box$ None of the above
\end{quote}


\end{enumerate}

\clearpage

\section{Recurrent Neural Networks (10 Points)}

\begin{enumerate}

\item Consider the following simple RNN architecture:

\begin{center}
\begin{tikzcd}[cells={nodes={draw=gray}}]
	& {\hat{y}_1} & {\hat{y}_2} & {\hat{y}_3} \\
	{h_0} & {h_1} & {h_2} & {h_3} \\
	& {x_1} & {x_2} & {x_3}
	\arrow[from=2-1, to=2-2]
	\arrow[from=2-2, to=2-3]
	\arrow[from=2-3, to=2-4]
	\arrow[from=3-4, to=2-4]
	\arrow[from=3-3, to=2-3]
	\arrow[from=3-2, to=2-2]
	\arrow[from=2-2, to=1-2]
	\arrow[from=2-3, to=1-3]
	\arrow[from=2-4, to=1-4]
\end{tikzcd}
\end{center}

where the dimensions of the layers and their corresponding weights are given below:
\begin{align*}
    \xv_t &\in \mathbb{R}^3 &\Wv_{hx} &\in \mathbb{R}^{4 \times 3}\\
    \hv_t &\in \mathbb{R}^4 &\Wv_{hy} &\in \mathbb{R}^{2 \times 4}\\
    \yv_t, \hat{\yv}_t &\in \mathbb{R}^2 &\Wv_{hh} &\in \mathbb{R}^{4 \times 4}
\end{align*}

The loss and update equations are:

\begin{align*}
    J &= \sum_{t=1}^3 J_t\\
    J_t &= -\sum_{i=1}^2 y_{t,i} \log(\hat{y}_{t,i})\\
    \hat{\yv}_t &= \sigma(\ov_t)\\
    \ov_t &= \Wv_{hy}\hv_t\\
    \hv_t &= \psi(\zv_t)\\
    \zv_t &= \Wv_{hh} \hv_{t-1} + \Wv_{hx}\xv_t\\
    \hv_0 &= 0
\end{align*}

Here, $\sigma$ is the \textbf{softmax} activation and $\psi$ is the \textbf{identity} activation (i.e., no activation). $J$ is the cross entropy loss and $t$ indicates timesteps. Note that we have no intercept terms (biases) in this architecture.

In this question, you will derive the steps of the backpropagation through time algorithm that lead to the computation of $\frac{dJ}{d\textbf{W}_{hh}}$. For all parts of this question, please write your answer in terms of $\Wv_{hh}$, $\Wv_{hy}$, $\yv_t$, $\hat{\yv}_t$, $\hv_t$, and any additional terms specified in the question. 

\clearpage

\begin{enumerate}
    \item \textbf{[2 pts]} What is $\Gv_{\ov_t} = \frac{\partial J}{\partial \ov_t}$? Write your solution in the first box, and show your work in the second. Write your answer in terms of $\hat{\yv}_t$ and $\yv_t$.
    

    \begin{your_solution}[title=$\frac{\partial J}{\partial \ov_t}$,height=2.5cm,width=10cm]
    % YOUR ANSWER

    \end{your_solution}
    
    \begin{your_solution}[title=Work,height=6cm,width=15cm]
    % YOUR ANSWER

    \end{your_solution}
    
    \item \textbf{[4 pts]} What is $\Gv_{\hv_i} = \frac{\partial J}{\partial \hv_i}$ for an arbitrary $i \in [1,3]$? Write your solution in terms of $\Gv_{\ov_t}$, $\Wv_{hh}$, $\Wv_{hy}$ in the first box, and show your work in the second. 
    
    \begin{your_solution}[title=$\frac{\partial J}{\partial \hv_i}$,height=2.5cm,width=8.5cm]
    % YOUR ANSWER

    \end{your_solution}
    
    \begin{your_solution}[title=Work,height=8cm,width=15cm]
    % YOUR ANSWER

    \end{your_solution}

    \clearpage
    
    \item \textbf{[4 pts]} What is $\Gv_{\Wv_{hh}} = \frac{\partial J}{\partial \Wv_{hh}}$? Write your solution in terms of $\Gv_{\hv_i}$ and $\hv$ in the first box, and show your work in the second.
    
    \begin{your_solution}[title=$\frac{\partial J}{\partial W_{hh}}$,height=2.5cm,width=10.5cm]
    % YOUR ANSWER

    \end{your_solution}
    
    \begin{your_solution}[title=Work,height=18cm,width=15cm]
    % YOUR ANSWER

    \end{your_solution}
\end{enumerate}

\end{enumerate}

\section{Attention \& Transformers (4 Points)}

\begin{enumerate}

\item \textbf{[2pts]} \textbf{Short Answer: } One limiting factor in a Transformer architecture for processing longer sequence lengths is its fixed context window. Does this limitation also apply to RNNs? Why or why not?

\begin{tcolorbox}[fit,height=4cm, width=15cm, blank, borderline={1pt}{-2pt},nobeforeafter]
%solution

\end{tcolorbox}


\item \textbf{[2 pts]} \textbf{Select One: } Given an input sequence of length $n$, how does the time complexity of the self-attention mechanism scale with $n$?
\begin{quote}
    $\circle$ $O(n$)\newline
    $\circle$ $O(n \log n$)\newline
    $\circle$ $O(n^2$)\newline
    $\circle$ $O(n^3$)
\end{quote}

\end{enumerate}